\begingroup
\raggedright
\subsection{Scientific Notation}
\texttt{aeb} and \texttt{aEb} mean $a\times10^b$. A positive
exponent shifts the decimal point right; a negative exponent shifts
it left. The exponent is an integer.

\begin{center}
\begin{tabular}{lll}
\textbf{Literal} & \textbf{Value} & \textbf{Type} \\
\hline
\texttt{1e3} & $1000$ & \texttt{double} \\
\texttt{2.5e-3} & $0.0025$ & \texttt{double} \\
\texttt{-4.2E+2} & $-420$ & \texttt{double} \\
\texttt{1e-9L} & $10^{-9}$ & \texttt{long double} \\
\texttt{1000000000LL} & $10^9$ (exact integer) & \texttt{long long} \\
\end{tabular}
\end{center}
Scientific literals are floating-point values, even \texttt{1e9}.
Use integer literals with \texttt{LL} when exact integer constants
are needed. Read exponent-form input with floating-point types or
\texttt{stod}/\texttt{stold}; integer parsers only consume the
initial integer prefix.

With \texttt{scientific} or \texttt{fixed}, \texttt{setprecision(p)}
sets digits after the decimal point. With \texttt{defaultfloat},
it sets significant digits. Stream flags and precision persist;
\texttt{defaultfloat} resets the float format, not the precision.
\texttt{uppercase}/\texttt{nouppercase} choose \texttt{E}/\texttt{e}.
Formatting rounds the display without changing the stored value.
\begin{samepage}
\medskip
\textbf{Quick fixed-format snippet.} Ten digits after the decimal point:
\cppfile{../more_topics/setprecision.cpp}
\end{samepage}
\lstinputlisting[style=C++,showstringspaces=false]{../compile/ScientificNotationExamples.cpp}
\endgroup
